%================================================
% RAPPORT PHASE 4 -- SUVA PROTOCOL
% SuVaWallet.sol -- Wallet ERC-4337 Double Signature
% Pour : William Schulz
% Par  : Mohamed Lamine MBENGUE
% Date : Avril 2026
%================================================
\documentclass[11pt,a4paper]{article}

\usepackage[utf8]{inputenc}
\usepackage[T1]{fontenc}
\usepackage[french]{babel}
\usepackage{geometry}
\geometry{margin=2.5cm}

\usepackage{amsmath,amssymb,amsthm}
\usepackage{booktabs}
\usepackage[table]{xcolor}
\usepackage{listings}
\usepackage{hyperref}
\usepackage{graphicx}
\usepackage{tikz}
\usetikzlibrary{arrows.meta, positioning, shapes.geometric, calc}
\usepackage{fancyhdr}
\usepackage{float}
\usepackage{array}
\usepackage{multirow}
\usepackage{tcolorbox}
\tcbuselibrary{skins, breakable}

%--- Couleurs ---
\definecolor{suva-blue}{RGB}{0, 82, 165}
\definecolor{suva-green}{RGB}{0, 150, 90}
\definecolor{suva-orange}{RGB}{220, 100, 0}
\definecolor{code-bg}{RGB}{245, 245, 245}
\definecolor{pass-green}{RGB}{0, 140, 70}
\definecolor{fail-red}{RGB}{180, 0, 0}

%--- En-tete/Pied de page ---
\pagestyle{fancy}
\fancyhf{}
\rhead{\textcolor{suva-blue}{\textbf{SuVa Protocol --- Rapport Phase 4}}}
\lhead{\textit{Confidentiel}}
\rfoot{\thepage}
\lfoot{\textit{Avril 2026}}

%--- Styles de code ---
\lstdefinestyle{solidity}{
  language=C,
  basicstyle=\ttfamily\small,
  backgroundcolor=\color{code-bg},
  keywordstyle=\color{suva-blue}\bfseries,
  commentstyle=\color{gray}\itshape,
  numbers=left, numberstyle=\tiny\color{gray},
  frame=single, breaklines=true, captionpos=b
}
\lstdefinestyle{bash}{
  language=bash,
  basicstyle=\ttfamily\small,
  backgroundcolor=\color{code-bg},
  keywordstyle=\color{suva-blue}\bfseries,
  commentstyle=\color{gray}\itshape,
  numbers=none,
  frame=single, breaklines=true, captionpos=b
}

%--- Boites personnalisees ---
\newtcolorbox{resultbox}[1][]{
  enhanced,colback=suva-green!8,colframe=suva-green!60,
  fonttitle=\bfseries,title=#1,arc=3mm
}
\newtcolorbox{warningbox}[1][]{
  enhanced,colback=suva-orange!8,colframe=suva-orange!60,
  fonttitle=\bfseries,title=#1,arc=3mm
}
\newtcolorbox{infobox}[1][]{
  enhanced,colback=suva-blue!6,colframe=suva-blue!50,
  fonttitle=\bfseries,title=#1,arc=3mm
}

%================================================
\begin{document}
%================================================

\begin{titlepage}
\centering
\vspace*{2cm}
{\Huge\bfseries\textcolor{suva-blue}{SuVa Protocol}\par}
\vspace{0.5cm}
{\Large\textit{Stablecoin Quantique sur EVM}\par}
\vspace{1.5cm}
\rule{\linewidth}{0.5pt}
\vspace{0.5cm}
{\LARGE\bfseries Rapport Technique --- Phase 4\par}
\vspace{0.3cm}
{\large SuVaWallet.sol --- Wallet ERC-4337 avec Validation Double Signature\par}
\vspace{0.5cm}
\rule{\linewidth}{0.5pt}
\vspace{2cm}

\begin{tabular}{ll}
\textbf{Client :}    & William Schulz \\[4pt]
\textbf{Auteur :}    & Mohamed Lamine MBENGUE \\[4pt]
\textbf{Portfolio :} & \href{https://portfolio-alamine.web.app}{\texttt{portfolio-alamine.web.app}} \\[4pt]
\textbf{Date :}      & Avril 2026 \\[4pt]
\textbf{Version :}   & 1.0 --- Final \\[4pt]
\textbf{Statut :}    & \textcolor{pass-green}{\textbf{COMPLET}} \\[4pt]
\textbf{Codebase :}  & ETHDILITHIUM-main (fork ZKNox) \\
\end{tabular}

\vfill
\begin{tcolorbox}[enhanced,colback=suva-blue!8,colframe=suva-blue,arc=4mm]
\centering\textit{Document confidentiel. Ce rapport couvre les livrables de la Phase 4 :\\
SuVaWallet.sol wallet ERC-4337, SuVaWalletFactory.sol, validation double signature\\
(ECDSA + Falcon-256), 15 tests forge, deploiement Sepolia.}
\end{tcolorbox}
\end{titlepage}

\tableofcontents
\newpage

%================================================
\section{Resume Executif}
%================================================

\begin{resultbox}[Phase 4 --- Tous les objectifs atteints]
\begin{itemize}
  \item[$\checkmark$] \textbf{SuVaWallet.sol} --- wallet ERC-4337 IAccount avec double validation
  \item[$\checkmark$] \textbf{SuVaWalletFactory.sol} --- factory CREATE2 pour deploiement deterministe
  \item[$\checkmark$] \textbf{Double signature} --- ECDSA + Falcon-256 toutes deux requises
  \item[$\checkmark$] \textbf{15 tests forge passants} --- tous les chemins de validation couverts
  \item[$\checkmark$] \textbf{41/41 tests totaux} sur toutes les phases --- zero regression
  \item[$\checkmark$] \textbf{Deploiement Sepolia} --- SuVaWallet live sur testnet
  \item[$\checkmark$] \textbf{Compatible EntryPoint ERC-4337 v0.6}
\end{itemize}
\end{resultbox}

\vspace{0.5cm}

La Phase 4 livre le \textbf{SuVaWallet} --- un smart contract wallet implementant le standard
ERC-4337 d'abstraction de compte. Chaque transaction doit etre autorisee par deux signatures
independantes : une signature ECDSA standard (compatible Ethereum) et une signature post-quantique
Falcon-256.

Cette architecture double signature offre une \textbf{defense en profondeur} : si les ordinateurs
quantiques cassent ECDSA a l'avenir, les fonds restent proteges par la signature Falcon. Si un bug
est trouve dans le verificateur Falcon, la signature ECDSA sert de repli.

%================================================
\section{Contexte ERC-4337}
%================================================

\subsection{Qu'est-ce que ERC-4337 ?}

ERC-4337 (Abstraction de Compte) permet aux smart contracts de fonctionner comme des wallets.
Au lieu de signer des transactions avec une cle privee, les utilisateurs soumettent des
\textbf{UserOperations} qui sont validees par la logique personnalisee de leur smart wallet.

\begin{center}
\begin{tikzpicture}[
  node distance=1.5cm,
  box/.style={rectangle,draw,rounded corners,minimum width=3cm,minimum height=0.9cm,
              align=center,font=\small},
  arrow/.style={-Stealth,thick}
]
  \node[box, fill=suva-blue!15]   (user)    at (0,0)    {Utilisateur};
  \node[box, fill=gray!15]        (bundler) at (4,0)    {Bundler};
  \node[box, fill=suva-orange!15] (ep)      at (8,0)    {EntryPoint};
  \node[box, fill=suva-green!15]  (wallet)  at (8,-2.2) {SuVaWallet};

  \draw[arrow] (user)    -- node[above,font=\footnotesize]{UserOperation} (bundler);
  \draw[arrow] (bundler) -- node[above,font=\footnotesize]{handleOps()} (ep);
  \draw[arrow] (ep)      -- node[right,font=\footnotesize]{validateUserOp()} (wallet);
  \draw[arrow] (wallet)  -- node[left, font=\footnotesize]{execute()} (ep);
\end{tikzpicture}
\end{center}

\begin{center}
\begin{tabular}{ll}
\toprule
\textbf{Acteur} & \textbf{Role} \\
\midrule
\textbf{Utilisateur} & Signe UserOperation avec ECDSA + Falcon \\
\textbf{Bundler}     & Collecte les UserOps et les soumet a l'EntryPoint \\
\textbf{EntryPoint}  & Contrat ERC-4337 officiel, valide et execute \\
\textbf{SuVaWallet}  & Logique de validation personnalisee : double signature \\
\bottomrule
\end{tabular}
\end{center}

\subsection{Pourquoi ERC-4337 pour SuVa ?}

Les wallets Ethereum standard utilisent une seule signature ECDSA. ERC-4337 permet a SuVa de :
\begin{itemize}
  \item Exiger \textbf{deux signatures} par transaction (ECDSA + Falcon)
  \item Definir une logique de validation personnalisee en Solidity
  \item Rester compatible avec l'ecosysteme ERC-4337 standard (bundlers, paymasters)
  \item Permettre des extensions futures (abstraction de gas, rotation de cle)
\end{itemize}

%================================================
\section{Architecture}
%================================================

\subsection{Vue d'ensemble des contrats}

\begin{center}
\begin{tabular}{lll}
\toprule
\textbf{Contrat} & \textbf{Role} & \textbf{Etat cle} \\
\midrule
\texttt{SuVaWallet.sol}        & Wallet ERC-4337, validation double sig & \texttt{owner}, \texttt{falconPkHash} \\
\texttt{SuVaWalletFactory.sol} & Factory CREATE2                        & sans etat \\
\texttt{IFalconVerifier.sol}   & Interface partagee (Phase 3 + 4)       & --- \\
\texttt{ZKNOX\_ethfalcon}      & Verificateur Falcon-256 (reutilise)    & sans etat \\
\bottomrule
\end{tabular}
\end{center}

\subsection{Encodage de la signature}

Le champ \texttt{userOp.signature} encode les deux signatures concatenees :

\begin{lstlisting}[style=bash, caption={Structure de la signature composite}]
userOp.signature = Sig_ECDSA (65 octets) || Payload_Falcon

ou :
  Sig_ECDSA      = abi.encodePacked(r, s, v)      -- 65 octets
  Payload_Falcon = abi.encode(h, salt, s1)
                    h    = uint256[256]  (cle publique Falcon)
                    salt = bytes         (nonce 40 octets)
                    s1   = uint256[256]  (polynome de signature)
\end{lstlisting}

\begin{infobox}[Decision de conception : signatures concatenees]
Concatener les deux signatures dans un seul champ \texttt{bytes} garde SuVaWallet compatible
avec tout bundler ERC-4337 sans modification. Le wallet lui-meme separe et valide chaque
signature independamment.
\end{infobox}

%================================================
\section{SuVaWallet.sol --- Implementation}
%================================================

\subsection{Variables d'etat}

\begin{lstlisting}[style=solidity, caption={Etat de SuVaWallet}]
IEntryPoint     public immutable entryPoint;
IFalconVerifier public immutable falconVerifier;

address public owner;        // proprietaire ECDSA
bytes32 public falconPkHash; // keccak256(abi.encodePacked(h))
\end{lstlisting}

Les deux adresses sont immuables --- definies a la construction, non modifiables.
\texttt{falconPkHash} stocke un engagement sur la cle publique Falcon (32 octets),
economisant 5,1M gas par rapport au stockage de la cle complete de 8 192 octets.

\subsection{validateUserOp()}

\begin{lstlisting}[style=solidity, caption={validateUserOp --- logique de double signature}]
function validateUserOp(
    UserOperation calldata userOp,
    bytes32               userOpHash,
    uint256               missingAccountFunds
) external override onlyEntryPoint returns (uint256 validationData) {
    // Payer le prefund EntryPoint si necessaire
    if (missingAccountFunds > 0) {
        (bool ok,) = payable(address(entryPoint)).call{
            value: missingAccountFunds
        }("");
        require(ok, "SuVaWallet: prefund failed");
    }

    // Rejeter si signature trop courte
    if (userOp.signature.length < 65) return SIG_VALIDATION_FAILED;

    // Separer la signature composite
    bytes memory ecdsaSig  = _slice(userOp.signature, 0, 65);
    bytes memory falconSig = _slice(userOp.signature, 65,
                                    userOp.signature.length - 65);

    // 1. Verifier ECDSA -- doit recuperer l'owner enregistre
    if (!_verifyECDSA(userOpHash, ecdsaSig))  return SIG_VALIDATION_FAILED;

    // 2. Verifier Falcon -- doit correspondre au hash de cle enregistre
    if (!_verifyFalcon(userOpHash, falconSig)) return SIG_VALIDATION_FAILED;

    return SIG_VALIDATION_SUCCESS; // == 0
}
\end{lstlisting}

\subsection{Verification ECDSA}

\begin{lstlisting}[style=solidity, caption={\_verifyECDSA}]
function _verifyECDSA(bytes32 hash, bytes memory sig)
    internal view returns (bool)
{
    if (sig.length != 65) return false;
    bytes32 r; bytes32 s; uint8 v;
    assembly {
        r := mload(add(sig, 32))
        s := mload(add(sig, 64))
        v := byte(0, mload(add(sig, 96)))
    }
    address recovered = ecrecover(hash, v, r, s);
    return recovered != address(0) && recovered == owner;
}
\end{lstlisting}

Utilise le precompile natif Ethereum \texttt{ecrecover}. L'adresse recuperee doit correspondre
exactement a \texttt{owner}.

\subsection{Verification Falcon}

\begin{lstlisting}[style=solidity, caption={\_verifyFalcon}]
function _verifyFalcon(bytes32 userOpHash, bytes memory falconSig)
    internal view returns (bool)
{
    if (falconPkHash == bytes32(0)) return false; // aucune cle enregistree

    (uint256[256] memory h, bytes memory salt, uint256[256] memory s1)
        = abi.decode(falconSig, (uint256[256], bytes, uint256[256]));

    // Verifier que la cle publique correspond au hash enregistre
    if (keccak256(abi.encodePacked(h)) != falconPkHash) return false;

    // Message signe = le userOpHash de 32 octets
    bytes memory message = abi.encodePacked(userOpHash);

    return falconVerifier.verify(h, salt, s1, message);
}
\end{lstlisting}

\begin{infobox}[Qu'est-ce que Falcon signe ?]
La signature Falcon est calculee sur le \texttt{userOpHash} --- un hash de 32 octets de toute
la UserOperation (sender, nonce, callData, limites de gas, etc.) calcule par l'EntryPoint.
Cela lie la signature Falcon a l'operation exacte autorisee, prevenant toute substitution
de parametre.
\end{infobox}

%================================================
\section{SuVaWalletFactory.sol}
%================================================

\begin{lstlisting}[style=solidity, caption={Factory --- deploiement deterministe CREATE2}]
function createWallet(
    address _entryPoint,
    address _falconVerifier,
    address _owner,
    uint256 salt
) external returns (SuVaWallet wallet) {
    bytes32 create2salt = keccak256(abi.encodePacked(_owner, salt));
    wallet = new SuVaWallet{salt: create2salt}(
        _entryPoint, _falconVerifier, _owner
    );
    emit WalletCreated(address(wallet), _owner);
}
\end{lstlisting}

\begin{infobox}[Pourquoi CREATE2 ?]
CREATE2 permet de calculer l'adresse d'un wallet \textit{avant} son deploiement, en utilisant
uniquement \texttt{owner} et \texttt{salt}. C'est requis par ERC-4337 : le bundler doit connaitre
l'adresse du wallet pour construire le champ \texttt{initCode} de la premiere UserOperation.
La fonction \texttt{getAddress()} retourne cette adresse predite.
\end{infobox}

%================================================
\section{Couverture des Tests}
%================================================

\subsection{Sortie forge test}

\begin{lstlisting}[style=bash, caption={forge test --match-path test/SuVaWallet.t.sol -vv}]
Ran 15 tests for test/SuVaWallet.t.sol:SuVaWalletTest
[PASS] testDeployment()                    (gas:    15 824)
[PASS] testRegisterFalconKey()             (gas:    69 526)
[PASS] testRegisterFalconKeyOnlyOwner()    (gas:    34 869)
[PASS] testValidateUserOpSuccess()         (gas: 5 322 396)
[PASS] testValidateUserOpWrongECDSA()      (gas: 5 198 564)
[PASS] testValidateUserOpWrongFalcon()     (gas: 5 323 328)
[PASS] testValidateUserOpShortSig()        (gas:    12 891)
[PASS] testValidateUserOpNoFalconKey()     (gas: 6 014 507)
[PASS] testValidateUserOpOnlyEntryPoint()  (gas: 1 147 169)
[PASS] testExecute()                       (gas:   107 194)
[PASS] testExecuteOnlyEntryPoint()         (gas:    79 643)
[PASS] testExecuteSendsETH()               (gas:    47 544)
[PASS] testFactoryCreateWallet()           (gas:   830 493)
[PASS] testFactoryGetAddress()             (gas:   833 663)
[PASS] testReceiveETH()                    (gas:    15 277)
Suite result: ok. 15 passed; 0 failed
\end{lstlisting}

\begin{resultbox}[15/15 tests Phase 4 passants --- 41/41 au total sur toutes les phases]
Zero echec. Zero regression.
\end{resultbox}

\subsection{Tableau des tests}

\begin{center}
\begin{tabular}{llc}
\toprule
\textbf{Test} & \textbf{Ce qu'il verifie} & \textbf{Resultat} \\
\midrule
\texttt{testDeployment}                   & owner, entryPoint, verifier corrects           & \textcolor{pass-green}{\textbf{OK}} \\
\texttt{testRegisterFalconKey}            & pkHash stocke comme keccak256(h)               & \textcolor{pass-green}{\textbf{OK}} \\
\texttt{testRegisterFalconKeyOnlyOwner}   & non-owner ne peut pas enregistrer              & \textcolor{pass-green}{\textbf{OK}} \\
\texttt{testValidateUserOpSuccess}        & ECDSA + Falcon valides $\Rightarrow$ retourne 0 & \textcolor{pass-green}{\textbf{OK}} \\
\texttt{testValidateUserOpWrongECDSA}     & mauvaise cle ECDSA $\Rightarrow$ retourne 1    & \textcolor{pass-green}{\textbf{OK}} \\
\texttt{testValidateUserOpWrongFalcon}    & sig Falcon invalide $\Rightarrow$ retourne 1   & \textcolor{pass-green}{\textbf{OK}} \\
\texttt{testValidateUserOpShortSig}       & sig $<$ 65 octets $\Rightarrow$ retourne 1     & \textcolor{pass-green}{\textbf{OK}} \\
\texttt{testValidateUserOpNoFalconKey}    & aucune cle enregistree $\Rightarrow$ retourne 1 & \textcolor{pass-green}{\textbf{OK}} \\
\texttt{testValidateUserOpOnlyEntryPoint} & appel non-EntryPoint revert                    & \textcolor{pass-green}{\textbf{OK}} \\
\texttt{testExecute}                      & appel transmis au contrat cible                & \textcolor{pass-green}{\textbf{OK}} \\
\texttt{testExecuteOnlyEntryPoint}        & execute non-EntryPoint revert                  & \textcolor{pass-green}{\textbf{OK}} \\
\texttt{testExecuteSendsETH}              & ETH transmis correctement                      & \textcolor{pass-green}{\textbf{OK}} \\
\texttt{testFactoryCreateWallet}          & factory deploie wallet avec bon owner          & \textcolor{pass-green}{\textbf{OK}} \\
\texttt{testFactoryGetAddress}            & adresse predite = adresse deployee             & \textcolor{pass-green}{\textbf{OK}} \\
\texttt{testReceiveETH}                   & wallet recoit ETH                              & \textcolor{pass-green}{\textbf{OK}} \\
\bottomrule
\end{tabular}
\end{center}

%================================================
\section{Deploiement Sepolia}
%================================================

Tous les contrats Phase 4 sont deployes et live sur le testnet Ethereum Sepolia (Chain ID : 11155111).

\begin{center}
\begin{tabular}{ll}
\toprule
\textbf{Contrat} & \textbf{Adresse Sepolia} \\
\midrule
\texttt{ZKNOX\_ethfalcon}      & \texttt{0xA6279C2CE813306da0f7B81E434E4A88f0b25626} \\
\texttt{MockERC20 (USDT)}      & \texttt{0xf088Bc3f822Ee138753AFe3190D1adEC60AbbaF3} \\
\texttt{qUSDT}                 & \texttt{0xf8f3009C473D538DA9B188F33bD6a6cBead29902} \\
\texttt{QuantumVault}          & \texttt{0x601CD837E9EDd0dfb311ff1A48eD7BeA7BBAf1c4} \\
\texttt{SuVaWalletFactory}     & \texttt{0x716a0C550eC1267Db493515d4a114C2FCf942e6B} \\
\texttt{SuVaWallet}            & \texttt{0x8a907C21014ED2926a53b4d18eC6e1b8C5a514cE} \\
\texttt{EntryPoint (ERC-4337)} & \texttt{0x5FF137D4b0FDCD49DcA30c7CF57E578a026d2789} \\
\bottomrule
\end{tabular}
\end{center}

\begin{center}
\begin{tabular}{ll}
\toprule
\textbf{Parametre} & \textbf{Valeur} \\
\midrule
Adresse deploieur & \texttt{0xD7390D91139Bf5868DDdF3b6677708b30b78D210} \\
Tx de deploiement & \texttt{0x73f4d8c1f997ca7620cf1a71af9cc5056ef8dcb56eeef4acc7b96fd5a800504e} \\
Bloc              & 10561066 \\
Gas total         & 4,790,419 gas \\
Cout total        & 0,000005748 ETH \\
Reseau            & Ethereum Sepolia (Chain ID 11155111) \\
\bottomrule
\end{tabular}
\end{center}

%================================================
\section{Decisions de Conception}
%================================================

\subsection{Les deux signatures sont-elles toujours requises ?}

\textbf{Decision : oui, les deux signatures requises pour chaque transaction.}

Rendre Falcon optionnel annulerait la resistance quantique --- un attaquant qui casse ECDSA
pourrait simplement soumettre des operations sans la signature Falcon. La friction UX
supplementaire est acceptable pour un wallet de garde de stablecoin ou la securite est
la priorite principale.

\subsection{La signature Falcon tient-elle dans les limites calldata ERC-4337 ?}

\begin{center}
\begin{tabular}{lrr}
\toprule
\textbf{Composant} & \textbf{Taille (octets)} & \textbf{Note} \\
\midrule
Signature ECDSA          &    65 & standard \\
\texttt{h} encode        & 8,192 & uint256[256] \\
\texttt{salt}            &    40 & octets nonce \\
\texttt{s1} encode       & 8,192 & uint256[256] \\
Overhead encodage ABI    &   128 & en-tetes abi.encode \\
\midrule
\textbf{Signature totale} & \textbf{16,617} & dans les limites bundler \\
\bottomrule
\end{tabular}
\end{center}

Les bundlers ERC-4337 autorisent typiquement jusqu'a 60 Ko de calldata. La signature
composite tient confortablement a \textasciitilde 16 Ko.

\subsection{IFalconVerifier refactored en interface partagee}

En Phase 4, \texttt{IFalconVerifier} a ete extrait de \texttt{QuantumVault.sol} vers un fichier
dedie \texttt{IFalconVerifier.sol}. Les deux contrats \texttt{QuantumVault} et \texttt{SuVaWallet}
importent depuis cette interface partagee, eliminant la duplication.

%================================================
\section{Analyse du Gas}
%================================================

\begin{center}
\begin{tabular}{llr}
\toprule
\textbf{Operation} & \textbf{Contrat} & \textbf{Gas} \\
\midrule
\texttt{validateUserOp} (MockFalcon)   & SuVaWallet & 5,322,396 \\
\texttt{validateUserOp} (mauvais ECDSA)& SuVaWallet & 5,198,564 \\
\texttt{validateUserOp} (mauvais Falcon)& SuVaWallet & 5,323,328 \\
\texttt{execute} (appel simple)        & SuVaWallet &   107,194 \\
\texttt{createWallet}                  & Factory    &   830,493 \\
\texttt{registerFalconKey}             & SuVaWallet &    69,526 \\
\bottomrule
\end{tabular}
\end{center}

%================================================
\section{Checklist des Livrables}
%================================================

\begin{center}
\begin{tabular}{clc}
\toprule
& \textbf{Livrable} & \textbf{Statut} \\
\midrule
$\checkmark$ & \texttt{SuVaWallet.sol} --- ERC-4337 IAccount          & \textcolor{pass-green}{\textbf{FAIT}} \\
$\checkmark$ & \texttt{SuVaWalletFactory.sol} --- factory CREATE2      & \textcolor{pass-green}{\textbf{FAIT}} \\
$\checkmark$ & Validation double signature (ECDSA + Falcon)            & \textcolor{pass-green}{\textbf{FAIT}} \\
$\checkmark$ & 15 tests forge passants                                  & \textcolor{pass-green}{\textbf{FAIT}} \\
$\checkmark$ & 41/41 tests totaux (toutes phases)                      & \textcolor{pass-green}{\textbf{FAIT}} \\
$\checkmark$ & Deploiement Sepolia                                      & \textcolor{pass-green}{\textbf{FAIT}} \\
$\checkmark$ & Analyse compatibilite bundler                            & \textcolor{pass-green}{\textbf{FAIT}} \\
$\checkmark$ & Questions de conception repondues                        & \textcolor{pass-green}{\textbf{FAIT}} \\
\bottomrule
\end{tabular}
\end{center}

\begin{resultbox}[Phase 4 : 8/8 objectifs completes]
Tous les livrables Phase 4 sont complets. SuVaWallet est live sur Sepolia.
\end{resultbox}

%================================================
\section{Conclusion}
%================================================

La Phase 4 livre un wallet quantique-resistant complet compatible ERC-4337. \textbf{SuVaWallet.sol}
impose une validation double signature --- chaque transaction requiert une signature ECDSA
(Ethereum standard) ET une signature post-quantique Falcon-256.

Le stack SuVa est maintenant complet jusqu'a la Phase 4 :

\begin{center}
\begin{tabular}{clcc}
\toprule
\textbf{Phase} & \textbf{Description} & \textbf{Statut} & \textbf{Tests} \\
\midrule
1  & Validation Dilithium / ETHDilithium       & \textcolor{pass-green}{\textbf{FAIT}} &  8/8 \\
2  & Integration Falcon / ETHFalcon            & \textcolor{pass-green}{\textbf{FAIT}} &  3/3 \\
3  & QuantumVault + qUSDT + Sepolia            & \textcolor{pass-green}{\textbf{FAIT}} & 15/15 \\
4  & SuVaWallet ERC-4337 + Sepolia             & \textcolor{pass-green}{\textbf{FAIT}} & 15/15 \\
5  & SPHINCS+ feasibilite (recherche)          & A venir & --- \\
\midrule
\multicolumn{2}{l}{\textbf{Total}} & & \textbf{41/41} \\
\bottomrule
\end{tabular}
\end{center}

%================================================
\appendix
\section{Annexe A --- Environnement}
%================================================

\begin{center}
\begin{tabular}{ll}
\toprule
\textbf{Outil} & \textbf{Version} \\
\midrule
Python          & 3.14 (environnement virtuel) \\
Foundry (forge) & Nightly (win32\_amd64) \\
Solidity        & 0.8.25 \\
Version EVM     & Prague \\
ERC-4337        & EntryPoint v0.6 \\
ZKNOX\_ethfalcon & scolaire (Phase 3/4) \\
\bottomrule
\end{tabular}
\end{center}

%================================================
\section{Annexe B --- Structure du Depot (Phase 4)}
%================================================

\begin{center}
\begin{tabular}{ll}
\toprule
\textbf{Chemin} & \textbf{Contenu} \\
\midrule
\texttt{src/SuVaWallet.sol}        & Wallet ERC-4337 double signature \\
\texttt{src/SuVaWalletFactory.sol} & Factory CREATE2 \\
\texttt{src/IFalconVerifier.sol}   & Interface verificateur Falcon partagee \\
\texttt{test/SuVaWallet.t.sol}     & 15 tests unitaires \\
\texttt{script/DeploySuVa.s.sol}   & Script deploiement complet (Phase 3 + 4) \\
\texttt{rapports/rapport\_phase4\_EN.tex} & Version anglaise \\
\texttt{rapports/rapport\_phase4\_FR.tex} & Ce document \\
\bottomrule
\end{tabular}
\end{center}

%================================================
\end{document}
%================================================
